Bevor die Anforderungen an das eigene System festgelegt werden, ordnet dieses Kapitel das Projekt in sein Umfeld ein.\kihinweis{}
Ausgangspunkt ist die Frage nach dem Publikum: wie Nachrichten heute auf mobilen Geräten genutzt werden und welche Angebote es für kurze Nachrichtenformate bereits gibt.
Danach folgen die zwei technischen Themen, auf denen wb-newsfeed aufbaut.
Das erste ist die automatische Zusammenfassung von Nachrichtentexten mit Sprachmodellen, das zweite die Personalisierung von Nachrichtenfeeds.
Zu beiden Themen gibt es umfangreiche Forschung, deren Ergebnisse mehrere Entscheidungen in diesem Projekt direkt beeinflusst haben.

\section{Nachrichtenkonsum auf mobilen Geräten}
\setauthor{\secondauthor}

Der jährlich erscheinende Digital News Report des Reuters Institute for the Study of Journalism beruht auf einer Online-Befragung in 48 Märkten und erlaubt dadurch Vergleiche über Länder und Jahre hinweg \cite{newman2025dnr}.
Die Ausgabe 2025 beschreibt eine deutliche Verschiebung hin zu sozialen Plattformen und Videoinhalten.
Der Anteil der Befragten, die mindestens einmal pro Woche Nachrichten als Social-Media-Video konsumieren, ist über alle Märkte von 52\,\% im Jahr 2020 auf 65\,\% im Jahr 2025 gestiegen.
Unter den 18- bis 24-Jährigen nennen 44\,\% soziale Netzwerke und Videoplattformen als ihre wichtigste Nachrichtenquelle \cite{newman2025dnr}.

Für Medienhäuser bedeutet diese Entwicklung, dass ein großer Teil des jungen Publikums Nachrichten zuerst in einem Feed sieht und nicht auf der eigenen Website.
Feeds auf Plattformen wie TikTok, Instagram Reels oder YouTube Shorts folgen einem gemeinsamen Muster: Ein Beitrag füllt den gesamten Bildschirm, der nächste erscheint nach einer Wischbewegung, und die Reihenfolge bestimmt ein Empfehlungsalgorithmus.
Ein Beitrag hat damit nur wenige Sekunden Zeit, um zu zeigen, worum es geht.

Ein weiterer Befund betrifft die Nachrichtenvermeidung.
Laut demselben Bericht geben 40\,\% der Befragten an, Nachrichten manchmal oder oft bewusst zu meiden, im Jahr 2017 waren es noch 29\,\% \cite{newman2025dnr}.
Der Bericht fragt außerdem, welche Angebote Nachrichten für die Befragten attraktiver machen würden.
27\,\% wünschen sich zusammengefasste Fassungen von Nachrichtenartikeln \cite{newman2025dnr}.
Zum ersten Mal wurde 2025 auch nach KI-Chatbots als Nachrichtenquelle gefragt.
Insgesamt nutzen 7\,\% solche Werkzeuge wöchentlich für Nachrichten, bei den unter 25-Jährigen sind es 15\,\% \cite{newman2025dnr}.

Aus diesen Zahlen lassen sich drei Folgerungen für einen Short-Form-Feed ableiten.
Erstens muss der Kern einer Nachricht ohne Öffnen des Artikels erkennbar sein, also in einer Überschrift, wenigen Sätzen und einem Bild.
Zweitens muss die Darstellung für Smartphones im Hochformat ausgelegt sein, weil dort die Nutzung stattfindet.
Drittens ist eine Kurzfassung kein Ersatz für den Artikel.
Wer sich für ein Thema interessiert, soll ohne Umweg zum vollständigen Text kommen.
Diese Folgerungen finden sich in den Anforderungen in Kapitel~\ref{chapter:requirements} wieder.
Offen bleibt, in welcher Form Medienhäuser solche Inhalte heute bereits anbieten, worum es im nächsten Abschnitt geht.

\section{Bestehende Short-Form-Angebote}
\setauthor{\secondauthor}

Kurze Nachrichtenformate sind nicht neu.
Die indische App Inshorts fasst Nachrichten aus verschiedenen Quellen in Texten mit höchstens 60 Wörtern zusammen und wirbt mit genau dieser Begrenzung \cite{inshorts}.
Das Angebot zeigt, dass ein Publikum für stark verkürzte Nachrichten existiert.
Es ist allerdings ein eigenständiges Produkt mit eigener Marke.
Ein Medienhaus, das seine Artikel dort erscheinen lassen möchte, hat keinen Einfluss darauf, wie die Kurzfassung entsteht, wie sie aussieht und welche Daten über die Nutzung gesammelt werden.

Der zweite verbreitete Weg ist die Veröffentlichung auf sozialen Plattformen.
Ein Medienhaus kann dafür eigene Kurzvideos oder Bildstrecken für Instagram und TikTok produzieren.
Damit erreichen sie das Publikum dort, wo es sich ohnehin aufhält, wie der vorige Abschnitt gezeigt hat.
Die Inhalte entstehen dabei aber in Handarbeit, meist zusätzlich zum eigentlichen Artikel.
Außerdem liegen Reichweite, Reihenfolge im Feed und Nutzungsdaten bei der Plattform und nicht beim Medium selbst.
Wer auf TikTok eine Kurzmeldung sieht, landet nicht automatisch auf der Website, von der die Meldung stammt.

Ein dritter Ansatz sind Feeds auf der eigenen Website oder in der eigenen App.
Hier behält das Medium die Kontrolle über Darstellung und Daten, muss aber Kurzfassungen und Bilder ebenfalls erst erstellen.
Für kleinere Redaktionen ist dieser zusätzliche Aufwand pro Artikel ein wesentliches Hindernis.

Tabelle~\ref{tab:related:comparison} stellt die drei Ansätze den Zielen von wb-newsfeed gegenüber.
Die Einträge der ersten drei Spalten beschreiben die Ansätze allgemein und nicht ein bestimmtes Produkt.

\begin{table}
    \centering
    \small
    \begin{tabularx}{\textwidth}{>{\raggedright\arraybackslash}X >{\raggedright\arraybackslash}X >{\raggedright\arraybackslash}X >{\raggedright\arraybackslash}X >{\raggedright\arraybackslash}X}
        \toprule
        \textbf{Kriterium} & \textbf{Aggregator-App} & \textbf{Soziale Plattform} & \textbf{Eigener Feed (manuell)} & \textbf{wb-newsfeed} \\
        \midrule
        Kurzfassung & extern erstellt & von Hand & von Hand & automatisch, von der Redaktion prüfbar \\
        Ort der Darstellung & fremde App & fremde Plattform & eigene Website & eigene Website \\
        Reihenfolge & Anbieter & Plattform-Algorithmus & Redaktion & Personalisierung im Browser \\
        Nutzungsdaten & beim Anbieter & bei der Plattform & beim Medium & beim Medium \\
        \bottomrule
    \end{tabularx}
    \caption{Vergleich verbreiteter Ansätze für Short-Form-Nachrichten}
    \label{tab:related:comparison}
\end{table}

Die Lücke, die das Projekt schließen soll, liegt in der Kombination aus der letzten Spalte.
Bestehende Artikel sollen ohne zusätzliche Handarbeit zu Feed-Inhalten werden, und der Feed soll als Teil der bestehenden Website laufen.
Damit das gelingt, muss die Kurzfassung automatisch entstehen.
Wie gut Sprachmodelle diese Aufgabe erfüllen und wo ihre Grenzen liegen, behandelt der folgende Abschnitt.

\section{Automatische Zusammenfassung mit Sprachmodellen}\label{sec:related:summarization}
\setauthor{\firstauthor}

In der Forschung zur automatischen Textzusammenfassung werden zwei Verfahren unterschieden.
Extraktive Verfahren wählen Sätze aus dem Originaltext aus und setzen sie zusammen.
Abstraktive Verfahren formulieren eine neue, kürzere Fassung in eigenen Worten \cite{maynez2020faithfulness}.
Für einen Teaser in einem Feed ist nur der zweite Weg sinnvoll, weil die ersten Sätze eines Nachrichtenartikels selten eine in sich geschlossene Kurzfassung ergeben.

Mit großen Sprachmodellen (\gls{llm}) hat sich die abstraktive Zusammenfassung stark verändert.
Goyal, Li und Durrett verglichen Zusammenfassungen von GPT-3, die ohne Training auf die Aufgabe nur durch eine Anweisung im Prompt entstanden, mit Modellen, die speziell für Nachrichtenzusammenfassung trainiert worden waren \cite{goyal2022news}.
Die befragten Personen bevorzugten in ihrer Studie deutlich die Texte von GPT-3.
Gleichzeitig zeigten die Autor:innen, dass gängige automatische Metriken diesen Unterschied nicht abbilden \cite{goyal2022news}.
Zhang et al.\ untersuchten zehn Sprachmodelle in einer großen Studie mit menschlicher Bewertung \cite{zhang2024benchmarking}.
Entscheidend für die Qualität war dabei nicht die Größe des Modells, sondern ob es auf das Befolgen von Anweisungen trainiert wurde.
Die Zusammenfassungen dieser Modelle wurden ähnlich gut bewertet wie Texte, die freiberufliche Autor:innen für die Studie geschrieben hatten \cite{zhang2024benchmarking}.
Für dieses Projekt ist der erste Befund besonders wichtig.
Er spricht dafür, dass auch ein mittelgroßes Modell, das lokal auf einer einzelnen Grafikkarte läuft, für Teaser ausreicht, solange es auf Anweisungen trainiert ist.

Das größte Problem abstraktiver Zusammenfassungen sind sogenannte Halluzinationen.
Maynez et al.\ zeigten an Nachrichtenzusammenfassungen, dass Modelle regelmäßig Inhalte erzeugen, die nicht durch den Ausgangstext gedeckt sind \cite{maynez2020faithfulness}.
Ji et al.\ unterscheiden in ihrer Übersichtsarbeit zwei Arten \cite{ji2023hallucination}.
Intrinsische Halluzinationen widersprechen dem Ausgangstext, etwa durch eine vertauschte Zahl.
Extrinsische Halluzinationen fügen Informationen hinzu, die sich aus dem Ausgangstext weder bestätigen noch widerlegen lassen.
Im Nachrichtenbereich sind beide Arten heikel, weil ein Teaser unter dem Namen der Redaktion erscheint.

Aus diesen Ergebnissen leiten sich mehrere Entscheidungen für wb-newsfeed ab, die in den folgenden Kapiteln genauer beschrieben werden.
Das Modell bekommt eine eng begrenzte Aufgabe mit festem Ausgabeformat: höchstens drei Sätze Teaser, drei Kernpunkte und Tags ausschließlich aus einer vorgegebenen Liste.
Titel und Kategorie, die die Redaktion bereits gesetzt hat, darf das Modell nicht überschreiben.
Jedes Ergebnis wird mit Modell und Parametern versioniert, damit sich nachvollziehen lässt, wie ein Text entstanden ist.
Die Redaktion kann jeden generierten Text im Admin-\gls{cms} prüfen und von Hand korrigieren.
Eingesetzt wird ein Modell der Qwen3-Familie, das sowohl einen Modus mit vorgeschaltetem Denkprozess als auch einen Modus für direkte Antworten bietet \cite{qwen2025qwen3}.
Für Teaser genügt der direkte Modus, für die anspruchsvollere Planung von Bildmotiven wird der Denkmodus genutzt.

Die Kurzfassung allein macht aber noch keinen Feed.
Ebenso wichtig ist, in welcher Reihenfolge die Beiträge erscheinen, was Thema des nächsten Abschnitts ist.

\section{Personalisierung von Nachrichtenfeeds}\label{sec:related:personalization}
\setauthor{\firstauthor}

Empfehlungssysteme für Nachrichten unterscheiden sich in einigen Punkten von Systemen für Filme oder Produkte.
Raza und Ding nennen in ihrer Übersichtsarbeit unter anderem die kurze Lebensdauer von Nachrichten, das Fehlen von Informationen über neue Nutzer:innen und neue Artikel sowie den Umstand, dass Interessen meist nur indirekt aus dem Verhalten ablesbar sind \cite{raza2022news}.
Ein Artikel über ein Fußballspiel ist am nächsten Tag kaum noch relevant, während ein Film jahrelang empfohlen werden kann.

Wu et al.\ veröffentlichten mit MIND einen großen Datensatz für Nachrichtenempfehlung, der aus den Klickdaten von Microsoft News entstanden ist und rund eine Million Nutzer:innen sowie über 160\,000 englischsprachige Artikel umfasst \cite{wu2020mind}.
Ihre Vergleiche bestehender Verfahren zeigen, dass die Qualität der Empfehlungen stark davon abhängt, wie gut ein System den Inhalt der Artikel versteht und wie gut es die Interessen der Nutzer:innen modelliert \cite{wu2020mind}.
Inhaltsbasierte Verfahren, die Artikel über ihren Text vergleichen, haben hier einen Vorteil.
Sie funktionieren auch für Artikel, die gerade erst veröffentlicht wurden und noch keine Klicks haben.

Eine verbreitete Technik für den inhaltlichen Vergleich sind \glspl{embedding}.
Ein Embedding-Modell bildet einen Text auf einen Vektor ab, und die Ähnlichkeit zweier Texte lässt sich dann über den Winkel zwischen ihren Vektoren bestimmen.
wb-newsfeed verwendet dafür das Modell BGE-M3, das über 100 Sprachen unterstützt und Texte mit bis zu 8\,192 Tokens verarbeiten kann \cite{chen2024bgem3}.
Für deutschsprachige Nachrichtenartikel ist die Mehrsprachigkeit eine Voraussetzung, die rein englischsprachige Modelle nicht erfüllen.

Personalisierung hat auch eine Kehrseite.
Pariser beschreibt unter dem Begriff Filterblase die Gefahr, dass Nutzer:innen durch personalisierte Auswahl nur noch Inhalte sehen, die ihre bisherigen Interessen bestätigen \cite{pariser2011filter}.
Ein bekanntes Gegenmittel aus dem Information Retrieval ist die Maximal Marginal Relevance von Carbonell und Goldstein \cite{carbonell1998mmr}.
Dabei wird bei der Sortierung nicht nur die Relevanz eines Dokuments berücksichtigt, sondern auch, wie sehr es sich von den bereits gewählten Dokumenten unterscheidet.

Für wb-newsfeed ergeben sich daraus zwei Leitlinien.
Die Personalisierung soll ohne Benutzerkonto funktionieren, weil ein eingebetteter Feed auf einer fremden Website keine Anmeldung voraussetzen kann.
Das Interessenprofil bleibt deshalb im Browser und wird nur aus dem Verhalten im Feed gelernt, etwa aus Verweildauer, Likes und gelesenen Artikeln.
Außerdem soll der Feed bewusst Abwechslung erzeugen.
Er begrenzt daher, wie viele Artikel desselben Themas hintereinander erscheinen, und reserviert regelmäßig Plätze für Themen außerhalb der bisherigen Interessen.
Wie diese Leitlinien in konkrete Anforderungen übersetzt wurden, beschreibt das nächste Kapitel.
